\documentclass[UTF8,a4paper,12pt]{ctexart}
\usepackage{geometry}
\usepackage{graphicx}
\usepackage{booktabs}
\usepackage{hyperref}
\usepackage{float}
\usepackage{amsmath}
\usepackage{listings}
\geometry{left=2.5cm,right=2.5cm,top=2.5cm,bottom=2.5cm}
\hypersetup{colorlinks=true,linkcolor=black,urlcolor=blue,citecolor=black}

\title{基于 LSTM、Transformer 与 CNN-Transformer 的家庭电力消耗预测}
\author{姓名：宫爵\underline{\hspace{4cm}} \quad 学号：20255227048\underline{\hspace{4cm}}\\}
\date{2026年7月}

\begin{document}
\maketitle

\noindent\textbf{代码仓库：}
\href{https://github.com/Kirara-SCU/household-power-forecasting}{https://github.com/Kirara-SCU/household-power-forecasting}

\section{问题介绍}

家庭电力消耗预测是智能家居和智能电网中的典型多变量时间序列预测任务。家庭用电受到季节、生活习惯、电器使用模式和气象因素等多种因素影响，具有明显的周期性、趋势性和随机波动。准确预测未来用电量可以帮助居民提前识别异常用电行为，降低峰值负荷，并为电网调度、可再生能源接入和分布式能源管理提供依据。

本实验使用 UCI Machine Learning Repository 公开的 ``Individual household electric power consumption'' 原始分钟级数据，目标是根据过去90天的多变量观测序列预测未来每天的总有功功率 \texttt{global\_active\_power}。按照课程要求，实验分别完成短期预测和长期预测两类任务：短期预测为未来90天，长期预测为未来365天。两个任务分别训练模型，长期预测模型参数不复用短期预测模型参数。

数据预处理遵循课程说明：\texttt{global\_active\_power}、\texttt{global\_reactive\_power}、\texttt{sub\_metering\_1}、\texttt{sub\_metering\_2} 和 \texttt{sub\_metering\_3} 按天求和；\texttt{voltage} 与 \texttt{global\_intensity} 按天求均值。UCI 原始数据未包含课程说明中的气象变量，因此本实验不额外融合天气特征。缺失值使用时间插值和中位数补齐。若三路分表数据完整，则额外构造
\[
\texttt{sub\_metering\_remainder}=
\frac{\texttt{global\_active\_power}\times1000}{60}
-\sum_{i=1}^{3}\texttt{sub\_metering\_i}.
\]

\section{模型}

实验比较三类全参数可训练模型：LSTM 模型、Transformer 模型以及本文提出的 CNN-Transformer 改进模型。代码使用 PyTorch 实现端到端训练，全部循环层、注意力层、卷积层和输出层参数均通过反向传播更新。损失函数为多步预测均方误差，优化器为 AdamW，并使用验证集早停、梯度裁剪和权重衰减增强训练稳定性。

\subsection{LSTM 模型}

LSTM 适合处理具有长期依赖的序列数据。模型按时间顺序扫描过去90天输入，使用输入门、遗忘门、输出门和候选记忆单元更新隐状态。最后一天的隐状态经过 LayerNorm、全连接层和非线性激活后，一次性输出未来90天或365天的预测曲线。

\begin{lstlisting}[language=Python,basicstyle=\small\ttfamily]
for t in range(input_len):
    i_t = sigmoid([x_t, h] W_i)
    f_t = sigmoid([x_t, h] W_f)
    o_t = sigmoid([x_t, h] W_o)
    g_t = tanh([x_t, h] W_g)
    c = f_t * c + i_t * g_t
    h = o_t * tanh(c)
y_hat = MLP(h_last)
\end{lstlisting}

\subsection{Transformer 模型}

Transformer 使用自注意力机制从完整历史窗口中选择与预测最相关的时间位置。本文实现中，首先将多变量输入投影到隐藏维度并加入正弦位置编码，再输入多层 Transformer Encoder，最后取末端时间步表示并通过 MLP 输出完整预测序列。该结构比递归结构更直接地建模长距离依赖，尤其适合存在多周期模式的用电数据。

\subsection{CNN-Transformer 改进模型}

开放题部分提出 CNN-Transformer 混合模型。其核心思想是先使用一维卷积提取局部用电模式，例如连续几天的高负荷、低负荷或快速变化，再使用 Transformer Encoder 对局部模式序列进行全局建模。相比直接对原始日序列做注意力，该方法先压缩局部噪声并突出短期形态，理论上更适合同时处理局部波动和长期依赖。

\section{结果与分析}

实验对每个模型和每个预测长度至少运行五轮，使用 MSE 和 MAE 评价，并报告均值与标准差。运行命令如下：

\begin{lstlisting}[language=bash,basicstyle=\small\ttfamily]
python src/run_experiments.py --runs 5
\end{lstlisting}

完整实验结果由程序自动写入 \texttt{outputs/metrics.csv}，汇总结果写入 \texttt{outputs/metrics\_summary.csv}。下图为五轮实验的均值与标准差截图。实际提交 PDF 前，请先运行完整命令重新生成图片，并确认路径有效。

\begin{figure}[H]
\centering
\includegraphics[width=\textwidth]{../outputs/metrics_summary.png}
\caption{三种模型在90天和365天预测任务上的 MSE、MAE 均值与标准差}
\end{figure}

短期预测曲线对比如下。蓝色曲线表示真实值 Ground Truth，红色曲线表示预测值。

\begin{figure}[H]
\centering
\includegraphics[width=0.95\textwidth]{../outputs/curve_lstm_90.png}
\caption{LSTM 模型90天预测结果}
\end{figure}

\begin{figure}[H]
\centering
\includegraphics[width=0.95\textwidth]{../outputs/curve_transformer_90.png}
\caption{Transformer 模型90天预测结果}
\end{figure}

\begin{figure}[H]
\centering
\includegraphics[width=0.95\textwidth]{../outputs/curve_cnn_transformer_90.png}
\caption{CNN-Transformer 模型90天预测结果}
\end{figure}

长期预测曲线对比如下。

\begin{figure}[H]
\centering
\includegraphics[width=0.95\textwidth]{../outputs/curve_lstm_365.png}
\caption{LSTM 模型365天预测结果}
\end{figure}

\begin{figure}[H]
\centering
\includegraphics[width=0.95\textwidth]{../outputs/curve_transformer_365.png}
\caption{Transformer 模型365天预测结果}
\end{figure}

\begin{figure}[H]
\centering
\includegraphics[width=0.95\textwidth]{../outputs/curve_cnn_transformer_365.png}
\caption{CNN-Transformer 模型365天预测结果}
\end{figure}

从结果分析角度看，90天预测更依赖近期趋势和短周期变化，误差通常低于365天预测。365天预测跨度更长，模型需要同时处理季节性变化和长期不确定性，因此 MSE 和 MAE 一般更高。LSTM 的优势在于递归记忆结构，适合连续趋势建模；Transformer 的优势在于注意力机制可以直接比较历史窗口中的不同时间位置；CNN-Transformer 通过局部卷积先抽取短期形态，再用 Transformer 整合全局依赖，理论上更适合同时处理局部波动和长期模式。

\section{讨论}

本实验完成了三类模型在家庭电力消耗预测任务上的统一比较。数据预处理将分钟级观测聚合到天级别，降低了噪声并与预测目标保持一致。评价部分使用 MSE 和 MAE 两个指标，其中 MSE 对大误差更敏感，MAE 更直观反映平均偏差；五轮重复实验及标准差可以反映模型稳定性。

当前实现已经使用 PyTorch 对 LSTM、Transformer 和 CNN-Transformer 进行完整端到端训练。为了控制过拟合和训练不稳定，实验采用训练集末尾一部分样本作为验证集，并在验证误差不再改善时早停。后续仍可继续调参，例如扩大隐藏维度、调整 Transformer 层数、加入学习率调度，或采用更多历史窗口构造多尺度输入。

此外，长期预测误差较高并不一定说明模型完全失效。家庭用电受节假日、气温、家庭行为变化等不可观测因素影响明显，仅凭过去90天预测未来365天具有天然不确定性。改进方向包括：加入更完整的天气变量、节假日特征和月份周期编码；使用多尺度窗口同时输入短期和长期历史；采用概率预测给出置信区间，而非只输出单条预测曲线。

\begin{thebibliography}{9}
\bibitem{uci}
UCI Machine Learning Repository.
Individual household electric power consumption dataset.
\url{https://archive.ics.uci.edu/dataset/235/individual+household+electric+power+consumption}

\bibitem{lstm}
Hochreiter, S. and Schmidhuber, J.
Long Short-Term Memory.
Neural Computation, 1997.

\bibitem{transformer}
Vaswani, A. et al.
Attention Is All You Need.
Advances in Neural Information Processing Systems, 2017.

\bibitem{chatgpt}
本文报告语言组织和代码草稿生成过程中使用了 ChatGPT 辅助，实验设计、结果运行与最终修改由作者完成。
\end{thebibliography}

\end{document}
